% Source: https://artofproblemsolving.com/community/c3h2350694_triangle_inequality_with_equilateral_triangle
% Stable identity: aops-topic-2350694-q01
\textbf{Three small equilateral triangles.} Draw through $P$ the lines parallel to the sides of the triangle. The lines through $P$ parallel to $AB$ and to $AC$ meet $BC$ at points $V$ and $U$. Triangle $PVU$ has its sides parallel to those of $ABC$, so it is equilateral, with apex $P$ and base $VU$ on $BC$; let $x$ be its side. In the same way we get an equilateral triangle with apex $P$ and base on $CA$, with side $y$, and one with base on $AB$, with side $z$ (dashed below).

\begin{asy}
usepackage("amsmath");
size(7cm);
defaultpen(fontsize(10pt));
pair B=(0,0), C=(1,0), A=(0.5,sqrt(3)/2), P=(0.42,0.24);
pair D=extension(A,P,B,C), E=extension(B,P,C,A), F=extension(C,P,A,B);
pair V=extension(P,P+(B-A),B,C), U=extension(P,P+(C-A),B,C);
pair Y1=extension(P,P+(A-B),C,A), Y2=extension(P,P+(C-B),C,A);
pair Z1=extension(P,P+(B-C),A,B), Z2=extension(P,P+(A-C),A,B);
fill(P--V--U--cycle,rgb(0.85,0.92,1)); fill(P--Y1--Y2--cycle,rgb(1,0.9,0.8)); fill(P--Z1--Z2--cycle,rgb(0.88,1,0.88));
draw(P--V--U--cycle,dashed+blue); draw(P--Y1--Y2--cycle,dashed+orange); draw(P--Z1--Z2--cycle,dashed+heavygreen);
draw(A--B--C--cycle,linewidth(0.9));
draw(A--D); draw(B--E); draw(C--F);
draw(P--D,red+linewidth(1)); draw(P--E,red+linewidth(1)); draw(P--F,red+linewidth(1));
dot(P); dot(D); dot(E); dot(F);
label("$B$",B,SW); label("$C$",C,SE); label("$P$",P+0.1*dir(-41));
label("$D$",D,S); label("$E$",E,NE); label("$F$",F,NW);
label("$V$",V,S,blue); label("$U$",U,S,blue);
label("$x$",(0.7*V+0.3*U)+(0,0.035),blue+fontsize(8pt));
label("$A$",A,N);
\end{asy}

\textbf{Step 1: $x+y+z=1$.} The height of triangle $PVU$ from $P$ is the distance $d_a$ from $P$ to $BC$, and an equilateral triangle of side $s$ has height $\dfrac{\sqrt3}{2}s$, so $x=\dfrac{2}{\sqrt3}d_a$. Similarly $y=\dfrac2{\sqrt3}d_b$ and $z=\dfrac2{\sqrt3}d_c$, where $d_b$, $d_c$ are the distances from $P$ to $CA$ and $AB$. The triangles $PBC$, $PCA$, $PAB$ fill triangle $ABC$, so
\[\dfrac12\left(d_a+d_b+d_c\right)\cdot1=[ABC]=\dfrac{\sqrt3}{4},\qquad d_a+d_b+d_c=\dfrac{\sqrt3}2,\]
and therefore $x+y+z=\dfrac2{\sqrt3}\cdot\dfrac{\sqrt3}{2}=1$.

\textbf{Step 2: $PD<x$.} Since $P$ is inside the triangle, $P=A+s(B-A)+t(C-A)$ with $s,t>0$. The ray from $P$ away from $A$, along line $AP$, has direction $P-A=s(B-A)+t(C-A)$. The rays $PV$ and $PU$ have directions $B-A$ and $C-A$ (they are parallel to $AB$ and $AC$ and point towards $BC$). A positive combination of the two directions points strictly inside the angle $VPU$, so this ray meets the segment $VU$ at a point strictly between $V$ and $U$. That point lies on line $AP$ and on $BC$, so it is $D$.

Now let $M$ be the midpoint of $VU$. Then $PM\perp VU$ and $MD<\dfrac x2$, so
\[PD^2=PM^2+MD^2<PM^2+\left(\dfrac x2\right)^2=PV^2=x^2,\]
that is, $PD<x$.

\textbf{Step 3.} In the same way $PE<y$ and $PF<z$. Adding,
\[PD+PE+PF<x+y+z=1.\qquad\blacksquare\]
